%ANS transactions template for use with ANS conference paper submission
%MASTER FOR JOURNAL PAPER
%DO NOT EDIT
\documentclass{anstrans}


\title{A \acl{cbs} case study for Nuclear Siting Tools}
\author{\textsuperscript{1,*}R. A. Borrelli $\sq$ \textsuperscript{2}Denia Djoki{\'c}}

\institute{
    \textsuperscript{1}\acl{ui} -- \acl{iff} $\sq$ 1776 Science Center Drive $\sq$ Idaho Falls\\
    \textsuperscript{2}\acl{um} -- \acl{fpz} $\sq$ 1109 Geddes Avenue $\sq$ Ann Arbor\\
    \textsuperscript{*}rborrelli@uidaho.edu; djokic@umich.edu\\
}


% Optional disclaimer - remove this command to hide
%\disclaimer{}

%%%%% general 
%packages and definitions
\usepackage{xargs} %for \newcommandx
\usepackage[inline]{enumitem}
\usepackage{nicefrac}
\usepackage{ulem}
\usepackage{graphicx} % allows inclusion of graphics
\graphicspath{{$TEXIMG}} %path to graphics
%\usepackage[pdftex,dvipsnames]{xcolor,colortbl} %change font color
%\usepackage[hidelinks]{hyperref}
%\hypersetup{draft}
%\usepackage{booktabs} % nice rules (thick lines) for tables
\usepackage{microtype} % improves typography for PDF
\usepackage[capitalise]{cleveref}
\usepackage{soul}
\usepackage{color}
%\usepackage[font=itshape]{quoting} %too much whitespace when page limits
\usepackage[switch,columnwise,pagewise]{lineno} %line numbering
\usepackage{longtable} %need to put label at top under caption then \\ - use spacing
%%%%%


%%%%% references
\usepackage[numbers,sort&compress]{natbib} %for [1-3] - nsf.bst, neup.bst
\setlength{\bibsep}{0pt} %sets space between references
%\renewcommand{\bibsection}{} %suppresses large 'references' heading
\renewcommand\bibpreamble{\vspace{-0.2\baselineskip}} %sets spacing after heading if not using default references heading
\makeatletter
\renewcommand{\@biblabel}[1]{#1.\hfill} %bibliography ordered list has numbers left flush
\makeatother
%%%%%


%%%%% watermark
%\usepackage[firstpage,vpos=0.63\paperheight]{draftwatermark}
%\SetWatermarkText{\shortstack{DRAFT\\do not distribute}}
%\SetWatermarkScale{0.20}
%%%%%


%%%%% toc and glossaries
\usepackage[acronym,nomain,nonumberlist]{glossaries}
%\makenoidxglossaries
%%%%%


%%%%% editing
\newcommand{\edit}[1]{\textcolor{blue}{#1}} %shortcut for changing font color on revised text
\newcommand{\fn}[1]{\footnote{#1}} %shortcut for footnote tag
\newcommand*\sq{\mathbin{\vcenter{\hbox{\rule{.3ex}{.3ex}}}}} %makes a small square as a separator $\sq$
%\newcommand{\sk}[1]{\sout{#1}} %shortcut for default strikethrough - do not sk through citep
\newcommand\sk{\bgroup\markoverwith{\textcolor{red}{\rule[0.5ex]{1pt}{1pt}}}\ULon} %strikethrough with red line; not in \section{}
\AtBeginEnvironment{quote}{\itshape}
%\st{} does strikethrough using soul package but does not like acronyms
%%%%%


%%%%% colors
%http://latexcolor.com/
%https://en.wikibooks.org/wiki/LaTeX/Colors#:~:text=black%2C%20blue%2C%20brown%2C%20cyan,be%20available%20on%20all%20systems.
%\definecolor{aliceblue}{rgb}{0.94, 0.97, 1.0}
%\definecolor{antiquewhite}{rgb}{0.98, 0.92, 0.84}
%\definecolor{lightmauve}{rgb}{0.86, 0.82, 1.0}
%\definecolor{brilliantlavender}{rgb}{0.96, 0.73, 1.0}
%\definecolor{brandeisblue}{rgb}{0.0, 0.44, 1.0}
%\definecolor{darkmidnightblue}{rgb}{0.0, 0.2, 0.4}
\newcommand{\x}{\cellcolor{aliceblue}} %use to shade in table cell
\newcommand{\y}{\cellcolor{lightgray}} %use to shade in table cell
\newcommand{\z}{\cellcolor{antiquewhite}} %use to shade in table cell
%%%%%


%%%%% acronyms
\newcommand{\acf}{\acrfull} %full acronym
\newcommand{\acl}{\acrlong} %long acronym
\newcommand{\acs}{\acrshort} %short acronym
\newcommand{\acfp}{\acrfullpl} %full acronym plural
\newcommand{\aclp}{\acrlongpl} %long acronym plural
\newcommand{\acsp}{\acrshortpl} %short acronym plural
\newcommand{\Acf}{\Acrfull} %full acronym first letter capital
\newcommand{\Acl}{\Acrlong} %long acronym first letter capital
%%%%%%


%%%%% todonotes
\setlength{\marginparwidth}{15mm}
\newcommandx{\cmt}[2][1=]{\todo[author=\textbf{STRUCTURE},tickmarkheight=0.15cm,linecolor=black,backgroundcolor=red!25,bordercolor=black,#1]{#2}}
\newcommandx{\con}[2][1=]{\todo[author=\textbf{CONTENT},tickmarkheight=0.15cm,linecolor=black,backgroundcolor=brilliantlavender,bordercolor=black,#1]{#2}}
\newcommandx{\rab}[2][1=]{\todo[noline,author=\textbf{RAB},backgroundcolor=Plum!25,bordercolor=black,#1]{#2}}
\newcommandx{\cas}[2][1=]{\todo[noline,author=\textbf{CASSIE},tickmarkheight=0.15cm,backgroundcolor=blue!25,bordercolor=black,#1]{#2}}
\newcommandx{\den}[2][1=]{\todo[author=\textbf{DENIA},tickmarkheight=0.15cm,backgroundcolor=OliveGreen!25,bordercolor=black,#1]{#2}}
%\newcommandx{\rab}[2][1=]{\todo[author=\textbf{RAB},tickmarkheight=0.15cm,linecolor=Plum,backgroundcolor=Plum!25,bordercolor=black,#1]{#2}}
%\newcommandx{\han}[2][1=]{\todo[author=\textbf{ATTN: Haney},tickmarkheight=0.15cm,linecolor=OliveGreen,backgroundcolor=OliveGreen!25,bordercolor=OliveGreen,#1]{#2}}
%\newcommandx{\jon}[2][1=]{\todo[author=\textbf{ATTN: Johnson},tickmarkheight=0.15cm,linecolor=blue,backgroundcolor=blue!25,bordercolor=blue,#1]{#2}}

%%% highlighting 
\DeclareRobustCommand{\hlf}[1]{{\sethlcolor{lightmauve}\hl{#1}}}
\makeatletter
    \if@todonotes@disabled
    \newcommand{\hlb}[2]{#1}
    \else
    \newcommand{\hlb}[2]{\rab{#2}\hlf{#1}}
    \fi
    \makeatother
%%%
%%%%%


%%%%% table alignments
\newcolumntype{L}[1]{>{\raggedright\let\newline\\\arraybackslash\hspace{0pt}}m{#1}} %uses \raggedright with m,p{} in table column
\newcolumntype{C}[1]{>{\centering\let\newline\\\arraybackslash\hspace{0pt}}m{#1}} %uses \raggedright with m,p{} in table column
\newcolumntype{R}[1]{>{\raggedleft\let\newline\\\arraybackslash\hspace{0pt}}m{#1}} %uses \raggedright with m,p{} in table column
%%%%%


%%%%% acronyms
\input{$ACRONYM/acronyms}
%%%%%


%%%%% linenumbering
%\linenumbers %toggle line numbers
%\pagewiselinenumbers %reset line numbers on new page
%\modulolinenumbers[5] %line numbering interval
%%%%%


\begin{document}


\section{Introduction} \label{sec-intro}
\subsection{\acl{snf} management in the \acl{usa}}
Nuclear energy expansion can can achieve strong energy security and independence, as well as economic growth for the \acs{usa} \cite{dod26a,atk25a,doe25a,lew25a,cat25a}. However, management of commercial \acf{snf} remains an unsolved, though not unsolvable, challenge \cite{ara26a,cfkoe25a,cfbor24a}. \acs{snf} is currently stored and managed in pools and on dry storage pads across thirty-three states at operating and decommissioned \acfp{npp} with a total of just over 90,000 \acs{mtu} \cite{gao25b}. The \acf{cgc} \cite{cfbor24a} received \acs{doe} support funding as one of ultimately twelve consortia from the 2021 \acl{caa} in which funds were provided for `\ldots nuclear waste disposal activities \ldots including interim storage activities\ldots' \cite{caa21a} for the purposes of building capacity for siting a \acf{cis} \cite{doe23c}. The \acs{cgc} will provide recommendations to \acs{doe} on process improvements. The \acl{cgc} is not asked to nor will designate a location or community for siting \cite{ara26a}.

\subsection{Motivation \& Goals}\label{sec-goals}
The 2024 Trump Administration has issued several executive orders focused on the accelerated growth of nuclear technologies. This includes a national target of 400 \acs{gwe} by 2050, deploying nuclear reactors at military bases, recycling and reprocessing, and increasing domestic fuel production \cite{doe25a}. However, additional growth requires a strong back end management policy. More facilities will be needed; e.g., storage, processing, fabrication, etc. Therefore, communities may want to understand the technical feasibility or capacity of a local site to support these capacities. To that end, there are some siting tools for nuclear facilities designed for a variety of outcomes that could be used by a community for site screening. To start, here, some of these tools are critically reviewed within the context of community accessibility and utility for site screening for a \acs{cis}. This paper is a subset of a larger study and features a selection of the tools. Use of the appropriate tool or tools could assist a community in developing a go/no-go protocol for site screening. The tools are analyzed with this approach in mind. 


\section{Regulations}
\subsection{Storage}
\acf{npp} siting is regulated under 10\acs{cfr}100. The \acs{nrc} issues operating licenses for independent storage; i.e., `away-from-reactor' of \acs{snf} under 10\acs{cfr}72. \acs{nrc} will perform a technical review in terms of safety and environmental impacts. An approved license is valid for 40 years. Subpart E contains siting evaluation factors. These are risk-informed, where the frequency and severity of external events, natural or man-made, must be assessed. 10\acs{cfr}103 contains additional site criteria specifically for dry cask storage. 

\subsection{Transport}\label{sec-transport}
\acs{snf} can be transported by truck, rail, or barge. Considerations include location of \acsp{npp}, distance to a \acs{cis}, existing infrastructure, etc. 10\acs{cfr}71 governs transport of any radioactive materials and includes seeking relevant security, safeguards, and security risk information from State, Tribal, and local officials. \acs{nrc} must review and approve a route plan \cite{nrc13a}. General criteria includes minimizing transit time, limiting travel through high population densities, avoiding bridges and tunnels, and traveling over advanced safety routes. Road route plans also stipulate identification of safe havens for temporary refuge or emergencies under 10\acs{cfr}73.37 coordinated with states \cite{nrc13a}. Transport routes will need to be formulated once a \acs{cis} is sited.


\section{Siting tools} \label{sec-tools}
\Cref{tab-tools-1} compiles a summary of the features and limits of the tools surveyed. Examples were simply chosen based on the areas where one of the author lives and not meant to suggest these would be possible sites. All figures are user generated. Screenshots are provided because the tools did not offer an option to download and image. 

\subsection{\acf{sta}}
\begin{figure}[htp]
    \centering
    \includegraphics[width=0.45\textwidth]{environmental-justice-tools/bonneville-stand-exploration.png}
    \caption{\acs{sta} example.}
    \label{fig-bonneville-stand}
\end{figure}
\acs{sta} is a web application to examine feasibility for advanced nuclear facility siting \cite{fpz23b}. \acs{sta} provides data on population, socioeconomics, safety, regulations, etc., while allowing comparative analysis for different sites. Though designed for advanced nuclear facilities, the focus is on \acs{npp} siting. Under the Site Exploration option, Bonneville County, Idaho, was selected. There are four main data layers; economic, reference, proximity, and safety. Brownfields, electric energy generators, navigable waterways, greater than 500 people per square mile, and protected lands were selected. \cref{fig-bonneville-stand} shows the large red areas are protected lands, the blue dots are electric energy generators, a single yellow dot west of Iona indicates a Brownfield site, and the lighter red is the high population density. Therefore, \cref{fig-bonneville-stand} shows locations in the county where a site might be located. \acs{sta} could be useful for a community if the screening criteria were developed for a \acs{cis}. A community would benefit most from the Site Exploration option. Information on regulatory data was not readily apparent and should be included. Additional information on the criteria would be beneficial. 

\subsection{\acs{tra}}\label{sec-tragis}
\begin{figure}[htp]
    \centering
    \includegraphics[width=0.45\textwidth]{environmental-justice-tools/tragis-rail.png}
    \caption{\acs{tra} example with population density.}
    \label{fig-tragis-rail}
\end{figure}
\begin{figure}[htp]
    \centering
    \includegraphics[width=0.45\textwidth]{environmental-justice-tools/tragis-naturita.png}
    \caption{\acs{tra} example for highway travel.}
    \label{fig-tragis-naturita}
\end{figure}
\acf{tra} is a \acs{gis} web-based tool developed by \acs{onl} for transportation routing. The tool includes databases for highway, rail, and waterways for the continental \acs{usa} and uses the LandScan population distribution model \cite{pet18b}. Users can select origin and destination, and criteria are then determined based on the options selected. Tribal lands on the routes are also identified. Transportation is optimized by minimizing the total impedance between the origin and destination, which includes overall distance, preferred routes, and travel time. Layers to select include county boundaries, critical infrastructure, population, etc. \cref{fig-tragis-rail} shows the rail network with population density of the \acs{usa}. \cref{fig-tragis-naturita} shows two highway routes from Columbia Generating Station to Naturita, CO, using highway travel. The southern route avoids Idaho, while the other includes Idaho. The Idaho route includes higher population centers. Tribal Lands are also included on the map in red. Data is generated for travel time, distance between nodes, and impedance. The orange dots are `nodes' that are included in the \acs{tra} database. Distance and time is given at each node. For the final node in Naturita, in avoiding Idaho, the distance is about 1300 miles and over 32 hours. Highway travel was selected because there is no node on rail to Naturita. A multitude of data layers can be selected; population density, critical infrastructure, highways, railways, etc. The list of nodes is extensive. However, the nodes are not shown on the map beforehand; they are listed on a dropdown menu after a state is selected for origin and destination. Multiple travel modes; e.g., rail and waterway, can be selected, but the destination nodes need to be known. The \acs{gui} is reasonably intuitive. A community might be able to use \acs{tra} to investigate locations of critical infrastructure or high populations and Tribal lands along a potential transportation route. The community itself would not be planning the transportation of \acs{snf} as part of developing an \acs{eoi}. However, \acs{tra} may aid in decisionmaking as to whether a community may want to avoid certain localities for transport. 

\subsection{\acl{lit}}
\begin{figure}[htp]
    \centering
    \includegraphics[width=0.45\textwidth]{environmental-justice-tools/lit-bonneville.png}
    \caption{\acs{lit} screening criteria for Bonneville, County, Idaho.}
    \label{fig-lit-bonneville}
\end{figure}
\begin{figure}[htp]
    \centering
    \includegraphics[width=0.45\textwidth]{environmental-justice-tools/lit-bonneville-assessment.png}
    \caption{\acs{lit} assessment for Bonneville, County, Idaho.}
    \label{fig-lit-bonneville-assessment}
\end{figure}
\begin{figure}[htp]
    \centering
    \includegraphics[width=0.45\textwidth]{environmental-justice-tools/lit-compare.png}
    \caption{\acs{lit} comparison with Bonneville County and southwest Wyoming.}
    \label{fig-lit-compare}
\end{figure}
The \acf{lit} tool was developed by \acf{onl}. \acs{lit} can be used to evaluate areas for storage \cite{le23a}. \acs{lit} can compare up to two locations for siting feasibility. There are three main areas of functionality -- Explore, to evaluate areas of interest; Access, to select assessment considerations; Compare, to create reports on the assessment considerations. A map of the \acs{usa} is presented first, where the user can select areas at the state or county level. Assessment considerations can then be determined, but the only option is `Initial Site Screening Criteria'. These are Critical Habitat for Endangered and Threatened Species, Coastal Barriers, and Federally Designated Lands. \cref{fig-lit-bonneville} shows the screening for Bonneville County, Idaho. The western part of the county does not have any criteria to screen, but the eastern part has 1--3 screening criteria. For Access, several considerations are provided; e.g., energy layers, floodplains, hazardous facilities, etc. \cref{fig-lit-bonneville-assessment} shows the map for Bonneville County with all the assessment considerations included. The western part of the county shows additional flagged land areas. The Compare component to the tool allows for the comparison of two regions for selected assessment considerations. Counties could not be selected; only highlighting an area on the map was available. The assessment considerations were the same. \cref{fig-lit-compare} shows the comparison to an area around Bonneville County with a similar size region of southwest Wyoming. Largely due to lower population density, Wyoming offers more land area. For a community interested in submitting an \acs{eoi}, \acs{lit} offers limited but useful capabilities. The lack of granularity to only the county level is limiting. There is a wide selection of assessment considerations. \acs{lit} perhaps could be used to examine screening criteria on a county level, and then compare smaller areas within the county in conjunction with another tool that could provide additional data at the county level.  

\subsection{\acl{str}}
\begin{figure}[htp]
    \centering
    \includegraphics[width=0.45\textwidth]{nuclear-fuel-cycle/back-end/nuclear-waste/start-columbia-population.png}
    \caption{\acs{str} example for \acs{snf} transport.}
    \label{fig-start-columbia-population}
\end{figure}
\acs{str} is a \acs{gis} decision support tool designed by \acs{doe} for planning \acs{snf} transport \cite{abk16a}. \acs{str} generates transportation routes by user selected options and filters, such as minimizing distance, population, etc. Infrastructure, police, airports, casinos, etc., can also be selected as sites to indicate on the transportation route. Options include selection of origin, destination, and transportation method \cite{con23a}. Different transportation modes can be selected for a single shipping route. \cref{fig-start-columbia-population} shows a route from Columbia Generating Station to a location outside of Pueblo using heavy haul truck, minimizing population with a buffer of 800 m for \acs{snf}. The route could also be filtered for travel time or distance or a combination. The total distance was 3500 miles. The route is rather circuitous due to the population constraint and the use of the heavy haul truck. The map does not provide a lot of detail. There is difficulty in selecting a city as a destination. While \acs{str} offers a wide variety of options for transportation analysis, a community could use the map to identify locations of critical infrastructure in the locality that may be desired to be avoided. 


\section{Discussion}
\subsection{Drawbacks}
There is some redundancy in the tools. \acs{tra} and \acs{str} are essentially the same, though \acs{tra} is easier to use and provides more detail. \acs{sta} and \acs{lit} both provide site screening, though \acs{sta} is more intuitive to use and provides more options. Given the redundancy, a community may not understand which tool could be optimal for their needs. A major drawback for a community is that these tools require registration. Communities may be wary of privacy issues or if the tool owners may be collecting data without consent. There is also a lack of portability in terms of downloading an image of the tool output for use in presentations, hearings, etc. Use of the tools was difficult to find; e.g., peer reviewed studies or reports. More transparency on data sources could create more confidence in selection. In short, for academic use, these tools offer interesting applications.  

\subsection{Outlook}
Both site screening and transportation could be useful to a community to aid in decisionmaking. It would be difficult to combine both components into a single tool. With additional customization, a screening tool in conjunction with a transportation tool can offer guidance for communities in decisionmaking for potential siting of a \acs{cis}, even if in terms of a go/no-go protocol. These tools reviewed here were fairly intuitive. Data sources and regulations are fairly complicated to understand. These existing tools may require additional explanation on their use for a community.  


%\setlength{\LTpre}{-0.03cm}
\setlength{\arrayrulewidth}{0.3mm}
\begin{scriptsize}
\begin{longtable}{|L{0.14\linewidth}|L{0.25\linewidth}|L{0.23\linewidth}|L{0.13\linewidth}|}
    \caption{Summary of tools}
        \label{tab-tools-1}\\
        \hline
        \textbf{Tool}
        &\textbf{Use}
        &\textbf{Features}
        &\textbf{Limits}
        \\
        \hline


        \href{https://fptz.org/}{\acs{sta} site}
        & For advanced reactor siting. 
        & Site Exploration methodology could be useful for community planners for screening. Selecting criteria was easy to perform.
        & Cannot download the map. Criteria lacked sufficient explanations. Registration required.
        \\
        \hline
        

        \href{https://webtragis.ornl.gov/login}{\acs{tra} site}
        & Modeling transportation routing. Uses \acs{gis} to identify optimal means of transportation for user selected parameters.
        & Intuitive \acs{gui} allows selection of a lot of layers including population density and Tribal lands. Can combine transport routes manually. 
        & Cannot download the map as a .jpg, etc. National parks, wilderness not included in layers. Cannot just click on the map to plan route; Transportation `nodes' are not shown on the map. Registration required.
        \\
        \hline


        \href{https://lite.ornl.gov/#/}{\acs{lit} site}
        & Evaluate areas of interest for interim storage using a pre-configured assessment considerations.
        & Presents a map of screening criteria and assessment considerations down to the county level. Areas can also be compared. Descriptions of the criteria and considerations is provided.  
        & Cannot download the map as a .jpg., etc. Lacking granularity. Registration required. 
        \\
        \hline

         
        \href{https://start.energy.gov/}{\acs{str} site}
        & Geospatial decision-support tool developed for evaluating routing options and other aspects of transporting \acs{snf}.
        & Variety of data layers. Printable map options. Can select different transportation modes for one shipping route. Can select for distance, time, population buffer. Can select a wide range of infrastructure to filter.
        & Hard to select cities, etc., for destinations. Have to zoom in significantly to select a destination. Map renders slowly. 
        \\
        \hline
\end{longtable}
\end{scriptsize}


\section{Summary remarks} \label{sec-sum}
Four tools have been currently reviewed in terms of community application in considering engaging in an \acl{eoi}. There are about a dozen more to analyze with most from government agencies. So far, the tools are not difficult to use but do have some limitations. It would not be difficult to optimize some of them for community use. For a community or group interested in an \acs{eoi}, \acs{sta} could be useful if modified for \acsp{cis}. \acs{tra} may offer more utility if transportation routes are of interest. \acs{lit} would be useful if the tool could scale down to the census tract level. Only showing county level screening is limiting because many western states have very large counties. For effective back-end fuel cycle management, community autonomy can be prioritized in part by providing the necessary tools for informed, independent decision-making.


\section{Acknowledgments} 
This work reflects technical work, funded by the \acl{doe} under Award Number DE-NE0009332 (DE-FOA-0002575). Funding does not reflect endorsement by \acs{doe}.


%\newpage



%%%%%%% citations
%
%%% make sure $BIBREF contains the path to each bib file in .bashrc
%
\bibliographystyle{ans}
\bibliography{
    araujo,
    borrelli,
    consent-siting,
    doe,
    dod,
    epa,
    environmental-justice-tools,
    federal,
    gao,
    jenkins-smith,
    learning,
    malin,
    nsf,
    nrc,
    nuclear-general,
    nuclear-power,
    nuclear-storage,
    nuclear-transport,
    nuclear-waste,
    nwtrb,
    participation-inclusion,
    socio-tech-economics,
    yucca-mountain
}
%
%%%%%%%


\end{document}
